\section{Preliminaries} \label{sec:preliminaries}

\subsection{Differentiable Support Function}
Throughout the paper, we use the union of differentiable support function (DSF) to represent the 3D shapes of the target shape and the objects in the inventory.
Given a compact convex body $\convex \subset \real^3$, the support function $\dsf_\convex :\real^3 \backslash\{0\} \to \real$ is defined as
\begin{equation*}
    \dsf_\convex(x) = \max_{s \in \convex} x\cdot s
\end{equation*}
where $\cdot$ denotes the standard inner product in $\real^3$.
The work of Lee et al. \cite{lee2023uncertain} presented a differentiable version of the support function $\dsf_\convex$ that represents a smooth shape $\convex$.

The function $\dsf_{\convex}(x)$ can be easily extended through rigid-body transformation as follows:
\begin{equation*}
    \dsf_{\bar{\convex}(q)}(x) = \dsf_\convex(R^\top x) + t \cdot x
\end{equation*}
for $q := (R,t) \in \SE(3)$ where $\bar{\convex}(q)=\{Rp + t\mid p \in \convex\}$ denotes a transformed $\convex$ under rigid-body transformation.
Given two DSF represented convex bodies $\convex_1(q_1)$ and $\convex_2(q_2)$ with SE(3) poses $q_1,q_2$, the signed distance $\sd(\convex_1(q_1), \convex_2(q_2))$ between the two convex bodies can be easily computed using Riemannian optimization techniques associated with $\dsf_{\convex_1(q_1)}$ and $\dsf_{\convex_2(q_2)}$ as presented in An, et al. \cite{an2024riemannianmtd}.
Moreover, due to the smoothness of DSFs, the signed distance $\sd(\convex_1(q_1), \convex_2(q_2))$ and the other contact information (e.g., normal vectors and contact points) are differentiable with respect to $q_1$ and $q_2$, which allows us to use gradient-based optimization techniques to optimize the placement poses of the objects in the later sections.

\subsection{Envelope Theorem}
Leveraged in many contact associated configuration optimization related works \cite{lee2025contactfactorgraph}, \cite{chen2025bodex}, \cite{zurbrugg2025graspqp}, the envelope theorem is a powerful tool that allows us to compute the gradient of an objective function with respect to a parameter when the objective function is defined as the optimal value of another optimization problem that depends on that parameter.
Consider an optimization problem as follows:
\begin{IEEEeqnarray*}{rCl}
    \min_{q,\lambda} &\quad& f(q,\lambda) \\
    \text{s.t. } &\quad& g_i(q,\lambda) \le 0,\,\, i=1,\cdots,m \\
    &\quad& h_j(q,\lambda) = 0 ,\,\, j=1,\cdots,n
\end{IEEEeqnarray*}
and the corresponding Lagrangian formulation:
\begin{equation*}
L(q,\lambda,\nu,\zeta) = f(q,\lambda) + \sum_i \nu_i g_i(q,\lambda) + \sum_j \zeta_j h_j(q,\lambda)
\end{equation*}
where $\nu = (\nu_1,\cdots,\nu_m)$ and $\zeta = (\zeta_1,\cdots,\zeta_n)$ are dual variables with $\nu_i \ge 0 $ for all $i=1,\cdots,m$.
For a fixed $\bar{q}$, let $\lambda^*(\bar{q})$, $\nu^*(\bar{q})$, $\zeta^*(\bar{q})$ be the optimal solution of the above optimization problem.
Then \textit{Envelop theorem} \cite{carter2001matheconomics} states that the derivative of $\bar{f}(q) := f(q,\lambda^*(q))$ can be computed as follows:
\begin{equation}
    \frac{d\bar{f}(q)}{dq} = \frac{\partial L(q,\lambda^*,\nu^*,\zeta^*)}{\partial q}. \label{eq:envelope}
\end{equation}

\subsection{Monte Carlo Tree Search (MCTS)}

Given a state space $\mathcal{S}$, an action space $\mathcal{A}$, a transition function $\transition:\mathcal{S}\times\mathcal{A}\to\mathcal{S}$, and a reward function $\reward:\mathcal{S}\times\mathcal{A}\to\real$, the Monte Carlo Tree Search (MCTS) is a popular algorithm for solving sequential decision-making problems that can be modeled as Markov Decision Processes (MDPs).

